\section{Conclusion}\label{sec:r34conclusion}
Neural Bellman Operators connect a learned scalar continuation to feasible
economic decisions. The connection is made at the level of decision-relevant
continuation errors rather than by interpreting a regression loss as a policy
loss. Full-action Bellman bounds then give policy comparisons whose domain,
continuation, and implementation requirements can be examined separately.

For the recursive capital economy, this connection is constructive. Economic
primitives determine a finite cold-training budget; a gradient--curvature
inequality identifies when approximate second-order information is sufficient;
and a noncommuting warm-start contraction admits bounded update errors and
changing continuation targets. A joint allocation pays for two-sided Gram
error, own-policy evaluation, and the actor solve. The resulting policy bound
holds against the full adapted finite-cost class rather than only the neural
or linear policies used to construct the candidate. Strict moment margins
also permit a finite-arithmetic realization with exact evaluation and
explicitly rounded factor updates.

The experiments retain a common economic acceptance rule and charge the
complete service. They distinguish a policy's certified accuracy from the
relative cost of constructing it. The structural solution in the quadratic
economy remains a legitimate comparator, and the adverse finite-law reuse
frontiers remain visible. Thus the positive training and policy results do
not depend on removing strong alternatives or relabeling an unfavorable
computation. The original controlled-diffusion model, recursive and strategic
applications, and their separate assumptions are preserved. Extending the
constructive conclusion to those other settings requires their own evaluation,
coverage, and transfer accounts; those accounts are not supplied merely by
using the same neural architecture.
